\chapter{Präpositionen}
\label{cha:Praepositionen}

Präpositionen verbinden Wörter und zeigen Beziehungen zwischen ihnen.

\section{Wechselpräpositionen (Dual Präpositionen)}
\label{sec:Wechselpraep}

Die Wechselpräpositionen können sowohl mit Dativ (Ort: Wo?) als auch mit Akkusativ (Richtung: Wohin?) verwendet werden:

\begin{center}
\begin{tabular}{|l|c|c|}
\hline
\textbf{Präposition} & \textbf{Wo? (Dativ)} & \textbf{Wohin? (Akkusativ)} \\ \hline
in & in der Schule & in die Schule \\ \hline
an & an der Wand & an die Wand \\ \hline
auf & auf dem Tisch & auf den Tisch \\ \hline
unter & unter dem Tisch & unter den Tisch \\ \hline
vor & vor dem Haus & vor das Haus \\ \hline
hinter & hinter dem Haus & hinter das Haus \\ \hline
neben & neben dem Haus & neben das Haus \\ \hline
zwischen & zwischen den Bäumen & zwischen die Bäume \\ \hline
über & über dem Tisch & über den Tisch \\ \hline
\end{tabular}
\end{center}

\begin{tcolorbox}[title=Wichtig]
\begin{itemize}
    \item \textbf{Ort (Stationär)}: Dativ
    \item \textbf{Richtung (Bewegung)}: Akkusativ
\end{itemize}
\end{tcolorbox}

\section{Übungen zu Wechselpräpositionen}
\label{sec:WechselpraepUebungen}

\begin{enumerate}
    \item Ich stelle die Blumen auf \underline{\hspace{1cm}} Tisch. (Richtung) \quad \textit{(Lösung: den)}
    \item Die Blumen stehen auf \underline{\hspace{1cm}} Tisch. (Ort) \quad \textit{(Lösung: dem)}
    \item Geh in \underline{\hspace{1cm}} Garten! (Richtung) \quad \textit{(Lösung: den)}
    \item Er ist in \underline{\hspace{1cm}} Garten. (Ort) \quad \textit{(Lösung: dem)}
    \item Häng das Bild an \underline{\hspace{1cm}} Wand. (Richtung) \quad \textit{(Lösung: die)}
    \item Das Bild hängt an \underline{\hspace{1cm}} Wand. (Ort) \quad \textit{(Lösung: der)}
    \item Setz dich vor \underline{\hspace{1cm}} Fernseher. (Richtung) \quad \textit{(Lösung: den)}
    \item Ich sitze vor \underline{\hspace{1cm}} Fernseher. (Ort) \quad \textit{(Lösung: dem)}
    \item Leg das Handy unter \underline{\hspace{1cm}} Kissen. (Richtung) \quad \textit{(Lösung: das)}
    \item Das Handy liegt unter \underline{\hspace{1cm}} Kissen. (Ort) \quad \textit{(Lösung: dem)}
    \item Wir laufen zwischen \underline{\hspace{1cm}} Bäume. (Richtung) \quad \textit{(Lösung: die)}
    \item Die Straße verläuft zwischen \underline{\hspace{1cm}} Bäumen. (Ort) \quad \textit{(Lösung: den)}
    \item Steck den Brief in \underline{\hspace{1cm}} Umschlag. (Richtung) \quad \textit{(Lösung: den)}
    \item Der Brief ist in \underline{\hspace{1cm}} Umschlag. (Ort) \quad \textit{(Lösung: dem)}
    \item Stell den Stuhl neben \underline{\hspace{1cm}} Sofa. (Richtung) \quad \textit{(Lösung: das)}
    \item Der Stuhl steht neben \underline{\hspace{1cm}} Sofa. (Ort) \quad \textit{(Lösung: dem)}
    \item Wir gehen über \underline{\hspace{1cm}} Brücke. (Richtung) \quad \textit{(Lösung: die)}
    \item Die Enten schwimmen unter \underline{\hspace{1cm}} Brücke. (Ort) \quad \textit{(Lösung: der)}
    \item Leg das Kleid über \underline{\hspace{1cm}} Stuhl. (Richtung) \quad \textit{(Lösung: den)}
    \item Das Kleid hängt über \underline{\hspace{1cm}} Stuhl. (Ort) \quad \textit{(Lösung: dem)}
\end{enumerate}

\section{Reine Dativ-Präpositionen}
\label{sec:DativPraep}

\begin{center}
\begin{tabular}{|l|l|}
\hline
\textbf{Präposition} & \textbf{Beispiel} \\ \hline
aus & Er kommt aus Berlin. \\ \hline
bei & Ich bin bei Maria. \\ \hline
mit & Ich fahre mit dem Bus. \\ \hline
nach & Nach dem Essen... \\ \hline
seit & Seit zwei Jahren... \\ \hline
von & Das Buch ist von ihm. \\ \hline
zu & Ich gehe zur Schule. \\ \hline
\end{tabular}
\end{center}

\section{Reine Akkusativ-Präpositionen}
\label{sec:AkkusativPraep}

\begin{center}
\begin{tabular}{|l|l|}
\hline
\textbf{Präposition} & \textbf{Beispiel} \\ \hline
bis & Bis morgen! \\ \hline
durch & Durch den Wald \\ \hline
gegen & Gegen die Wand \\ \hline
ohne & Ohne dich \\ \hline
um & Um die Ecke \\ \hline
\end{tabular}
\end{center}

\section{Genitiv-Präpositionen}
\label{sec:GenitivPraep}

\begin{center}
\begin{tabular}{|l|l|}
\hline
\textbf{Präposition} & \textbf{Beispiel} \\ \hline
trotz & Trotz des Regens... \\ \hline
wegen & Wegen des Verkehrs... \\ \hline
während & Während des Films... \\ \hline
statt & Statt des Bruders... \\ \hline
außerhalb & Außerhalb der Stadt \\ \hline
innerhalb & Innerhalb der Woche \\ \hline
aufgrund & Aufgrund der Situation... \\ \hline
oberhalb & Oberhalb des Berges \\ \hline
unterhalb & Unterhalb des Tales \\ \hline
diesseits & Diesseits des Flusses \\ \hline
jenseits & Jenseits des Meeres \\ \hline
\end{tabular}
\end{center}

\section{Übungen zu Genitiv-Präpositionen}
\label{sec:GenitivPraepUebungen}

\begin{enumerate}
    \item \underline{\hspace{1cm}} schlechten Wetters blieben wir zu Hause. (trotz) \quad \textit{(Lösung: Trotz des)}
    \item \underline{\hspace{1cm}} seiner Krankheit fehlte er. (wegen) \quad \textit{(Lösung: Wegen)}
    \item \underline{\hspace{1cm}} des Films schlief ich ein. (während) \quad \textit{(Lösung: Während)}
    \item \underline{\hspace{1cm}} des Verkehrs kam er zu spät. (wegen) \quad \textit{(Lösung: Wegen)}
    \item \underline{\hspace{1cm}} des Kaffees trinke ich Tee. (anstatt) \quad \textit{(Lösung: Anstatt)}
    \item \underline{\hspace{1cm}} seiner Krankheit blieb er zu Hause. (aufgrund) \quad \textit{(Lösung: Aufgrund)}
    \item \underline{\hspace{1cm}} des Parks ist es verboten. (innerhalb) \quad \textit{(Lösung: Innerhalb)}
    \item \underline{\hspace{1cm}} der Öffnungszeiten ist geschlossen. (außerhalb) \quad \textit{(Lösung: Außerhalb)}
    \item \underline{\hspace{1cm}} der Baumgrenze wächst nichts. (oberhalb) \quad \textit{(Lösung: Oberhalb)}
    \item \underline{\hspace{1cm}} des Meeres ist Europa. (diesseits) \quad \textit{(Lösung: Diesseits)}
    \item \underline{\hspace{1cm}} der Berge liegt ein Tal. (jenseits) \quad \textit{(Lösung: Jenseits)}
\end{enumerate}

\section{Feste Präpositionen mit festem Kasus}
\label{sec:FestePraep}

\begin{tcolorbox}[title={Präposition: Dativ}]
aus, bei, mit, nach, seit, von, zu
\end{tcolorbox}

\begin{tcolorbox}[title={Präposition: Akkusativ}]
bis, durch, gegen, ohne, um
\end{tcolorbox}

\begin{tcolorbox}[title={Präposition: Genitiv}]
trotz, wegen, während, statt
\end{tcolorbox}

\chapterend